%!TEX root = ../dissertation.tex

Container terminals use simulation to test operational changes without
disrupting live operations, but these models are usually built and calibrated
by hand. Data-aware process simulation aims to learn executable models from
event logs. Evidence from physical logistics systems is still limited.

This thesis evaluates the approach for truck operations at HHLA Container
Terminal Burchardkai. Two months of data were transformed into an event log
with 89,460 truck visits and split temporally into training and test periods.
The white-box ProSiT framework was used to compare a distribution-only model
with a model using decision rules and workload features. The first
Inductive-Miner import appeared to allow impossible parallel behaviour within
one truck visit because minute-resolved timestamps were used to reorder events.
Preserving the explicit XES event order instead yields a compact sequential
Petri net, which is used unchanged as the common executable control flow.

The Petri net fits 99.994\,\% of the 17,892 held-out visits and contains no
parallel operator. Its loop generalises beyond exact training variants, while
a rarely sampled silent gate-only path is retained and disclosed as model
overgeneralisation rather than removed by an artificial constraint. The final
model reproduced arrival spacing closely and parts of yard service times, but
mean truck turnaround was 8.17 minutes too short and the 90th percentile was
17.31 minutes too short. After a common calibration, contextual rules and
workload features did not improve fidelity uniformly: timing distances remain
similar, the activity-rate error increases, and the rare silent bypass is no
longer sampled. RMG capacity still required a documented domain correction.
Two what-if experiments were transparent and repeatable, but their
mean-turnaround confidence intervals included zero. The model also lacked the
spatial state needed for physical counterfactual claims.

Data-aware process simulation can provide an auditable historical model and a
useful diagnostic tool for container-terminal operations. Direct operational
decision support requires explicit domain validation and richer models of
location, resources, and queueing. With these mechanisms validated, the
approach can support the central logistics use case: comparing candidate truck
slot rates and handling policies under forecast terminal conditions before
they are introduced in live operations.
